\documentclass[../../../include/open-logic-section]{subfiles}

\begin{document}

\olfileid{sth}{ordinals}{wo}
\olsection{في الترتيبات الجيدة}

والحد الأساس هو:

\begin{defn}
تُسمّى العلاقة~$<$ \emph{ترتيبًا جيدًا} على~$A$ إذا وفقط إذا حققت
الشرطين الآتيين:
\begin{enumerate}
\item أن تكون~$<$ متصلة: فلكل~$a, b \in A$ إما~$a < b$،
	أو~$a = b$، أو~$b < a$؛
\item وأن يكون لكل مجموعة جزئية غير خالية من~$A$ !!{element} أصغري بحسب~$<$،
	أي إذا كان~$\emptyset \neq X \subseteq A$، فإن
	$(\exists m \in X)(\forall z \in X)z \nless m$.
\end{enumerate}
\end{defn}

ومن اليسير التحقق من أن علاقات الأمثلة الثلاثة السابقة ترتيبات جيدة.

\begin{prob}
في~\Olref[sth][ordinals][idea]{sec} ثلاثة ترتيبات نموذجية على الأعداد
الطبيعية؛ أثبت أن كل واحد منها ترتيب جيد.
\end{prob}

وفي الترتيب الجيد ملاحظات أولية عظيمة الأثر:

\begin{prop}\ollabel{wo:strictorder}
إذا كانت~$<$ ترتيبًا جيدًا على~$A$، كان لكل مجموعة جزئية غير خالية من~$A$
عنصر أصغر وحيد بحسب~$<$، وكانت~$<$ غير انعكاسية ولا تناظرية ومتعدية.
\end{prop}

\begin{proof}
لتكن~$X$ مجموعة جزئية غير خالية من~$A$. فلها !!{element} أصغري بحسب
$<$، سمّه~$m$؛ أي $(\forall z \in X)z \nless m$. ومن اتصال~$<$ يلزم
$(\forall z \in X)m \leq z$؛ فيكون~$m$ أصغر !!{element} بحسب~$<$
في~$X$.

  ولعدم الانعكاس ثبّت~$a \in A$؛ فإن أصغر !!{element} بحسب~$<$ في
  $\{a\}$ هو~$a$، ومن ثم~$a \nless a$. ولعدم التناظر، لو كان
  $a<b$ و$b<a$ لما كان واحد من~$a$ و$b$ أصغريًا في~$\{a,b\}$، وذلك
  محال. وللتعدي افرض~$a<b<c$، وليكن~$m$ أصغر !!{element} في
  $\{a,b,c\}$. لا يكون~$m=b$ لأن~$a<b$، ولا يكون~$m=c$ لأن~$b<c$؛
  فلا بد أن يكون~$m=a$، ومن ثم~$a\leq c$. ولما كان~$a\neq c$ بحكم
  عدم التناظر، لزم~$a<c$، فثبت أن~$<$ متعدية.
\end{proof}

\begin{prop}\ollabel{propwoinduction}
إذا كانت~$<$ ترتيبًا جيدًا على~$A$، كان، لكل صيغة~$\phi(x)$:
\[
	\text{إذا كان }(\forall a \in A)((\forall b < a)\phi(b) \lif 
		\phi(a))\text{، فإن }(\forall a \in A)\phi(a).
\]
\end{prop}

\begin{proof}
نبرهن نقيض المقابل. افترض~$\lnot(\forall a \in A)\phi(a)$؛ أي إن
$X = \Setabs{x \in A}{\lnot\phi(x)} \neq \emptyset$. عندئذ يكون
لـ$X$ !!{element} أصغري بحسب~$<$، سمّه~$a$. ومن ثم
$(\forall b < a)\phi(b)$، مع أن~$\lnot \phi(a)$.
\end{proof}\noindent
وهذه الخاصية الأخيرة نظيرة مبدأ الاستقراء القوي على الأعداد الطبيعية،
وهو: إذا كان
$(\forall n \in \omega)((\forall m < n)\phi(m) \lif \phi(n))$، فإن
$(\forall n \in \omega)\phi(n)$. وبها يغدو الترتيب الجيد مفهومًا بالغ
\emph{القوة}.
\footnote{تذكير: يجوز أن تكون لجميع الصيغ معاملات (ما لم يُصرح بخلاف ذلك).}

\end{document}

